\documentclass[11pt, a4paper]{article}

% ── Page layout ───────────────────────────────────────────────────────────────
\usepackage[top=0.9in, bottom=0.9in, left=1.1in, right=1.1in]{geometry}

% ── Typography ────────────────────────────────────────────────────────────────
\usepackage{mathptmx}
\usepackage{microtype}
\usepackage{parskip}
\setlength{\parskip}{5pt}

% ── Mathematics ───────────────────────────────────────────────────────────────
\usepackage{amsmath}
\usepackage{amssymb}

% ── Figures and tables ────────────────────────────────────────────────────────
\usepackage{graphicx}
\usepackage{booktabs}
\usepackage{caption}
\captionsetup{font=small, labelfont=bf}

% ── Section formatting ────────────────────────────────────────────────────────
\usepackage{titlesec}
\titleformat{\section}{\large\bfseries}{\thesection.}{0.6em}{}
\titleformat{\subsection}{\normalsize\bfseries}{\thesubsection.}{0.5em}{}
\titlespacing*{\section}{0pt}{12pt}{5pt}
\titlespacing*{\subsection}{0pt}{8pt}{3pt}

% ── Hyperlinks ────────────────────────────────────────────────────────────────
\usepackage[colorlinks=true, linkcolor=black, citecolor=black, urlcolor=blue]{hyperref}

% ── Miscellaneous ─────────────────────────────────────────────────────────────
\usepackage{enumitem}
\setlist[itemize]{noitemsep, topsep=2pt}
\setlist[enumerate]{noitemsep, topsep=2pt}
\usepackage{xcolor}
\usepackage{fancyhdr}

\pagestyle{fancy}
\fancyhf{}
\rhead{\small EE559 --- Final Project Report}
\lhead{\small Om Suresh Prajapati}
\cfoot{\thepage}
\renewcommand{\headrulewidth}{0.4pt}

% ─────────────────────────────────────────────────────────────────────────────
\begin{document}

% ── Title block ──────────────────────────────────────────────────────────────
\begin{center}
    {\LARGE \textbf{Adaptive Fraud Detection with Dynamic Thresholding\\
    and Cost-Sensitive Learning}}\\[0.5em]
    {\large \textbf{Om Suresh Prajapati}} \quad|\quad USC ID: 1430823821\\[0.15em]
    \textit{EE559 (CSCI 559) --- Machine Learning, Spring 2026}\\[0.15em]
    \textit{Instructor: Prof.\ Anand A.\ Joshi, University of Southern California}\\[0.15em]
    May 4, 2026
\end{center}

\vspace{0.3em}
\hrule
\vspace{0.5em}

% ── Abstract ─────────────────────────────────────────────────────────────────
\noindent\textbf{Abstract} --- We present a complete supervised ML pipeline for credit
card fraud detection on the Kaggle Credit Card Fraud dataset (284,807 transactions;
0.17\% fraud prevalence). Three classifiers are trained and compared: Logistic
Regression (LR), Random Forest (RF), and Histogram-based Gradient Boosted Trees (GBT).
Three class-imbalance strategies are evaluated: balanced class weighting, SMOTE, and
random under-sampling. We introduce dynamic threshold tuning on the validation set and
a cost-sensitive evaluation framework (FN\,=\,\euro500 vs.\ FP\,=\,\euro10). The best
model---RF with balanced class weights at a tuned threshold of 0.192---achieves
\textbf{PR-AUC\,=\,0.8116}, \textbf{Recall\,=\,82.4\%}, and
\textbf{Precision\,=\,84.7\%} on the held-out test set.

\vspace{0.3em}
\hrule
\vspace{0.5em}

% ─────────────────────────────────────────────────────────────────────────────
\section{Introduction}

Credit card fraud causes billions of euros in annual losses. Automated detection
reduces this, but the problem presents two fundamental challenges. First, the
extreme class imbalance (577:1 in this dataset) renders standard accuracy
uninformative: labelling every transaction as legitimate achieves 99.83\% accuracy
while catching zero fraud. Second, the two error types carry asymmetric costs:
a missed fraud (false negative, FN) costs far more than a false alarm (false
positive, FP) that triggers a manual review.

This work addresses both challenges through appropriate primary metrics
(PR-AUC and ROC-AUC), three imbalance-handling strategies, and a cost-sensitive
threshold selection procedure that minimises expected business cost on the
validation set.

% ─────────────────────────────────────────────────────────────────────────────
\section{Dataset and Preprocessing}

\subsection{Dataset}

The Kaggle Credit Card Fraud Detection dataset~\cite{kaggle_ccfd} contains 284,807
European cardholder transactions (September 2013), with 492 fraud cases (0.173\%).
Features \texttt{V1}--\texttt{V28} are anonymised PCA components (zero-mean,
decorrelated); \texttt{Time} and \texttt{Amount} are in original units.
No missing values; 1,081 duplicate rows retained.

\subsection{Preprocessing Pipeline}

The following strict ordering prevents data leakage:
\begin{enumerate}
    \item \textbf{Split first.} Stratified 70/15/15 train/validation/test split
          (199,356 / 42,729 / 42,722 samples), preserving 0.173\% fraud in each subset.
    \item \textbf{Scale on train only.} \texttt{StandardScaler} fitted on
          \texttt{Time} and \texttt{Amount} from training data; applied to all splits.
          Features \texttt{V1}--\texttt{V28} require no re-scaling (PCA output).
    \item \textbf{Resample train only.} SMOTE and under-sampling applied within the
          training split only, never touching validation or test.
\end{enumerate}

% ─────────────────────────────────────────────────────────────────────────────
\section{Methodology}

\subsection{Models}

\textbf{Logistic Regression (LR).}
Linear baseline optimising cross-entropy with L2 regularisation ($C$\,=\,1.0).
\texttt{class\_weight=`balanced'} reweights each fraud sample by
$\approx$577$\times$ relative to legitimate samples.

\textbf{Random Forest (RF).}
Ensemble of 100 trees trained on bootstrap samples with random feature subsets.
Captures nonlinear feature interactions; same balanced class weighting.

\textbf{Gradient Boosted Trees (GBT).}
\texttt{HistGradientBoostingClassifier} (sklearn, 200 iterations). Additively builds
trees to correct ensemble residuals. Two variants: balanced class weighting (GBT-bal)
and SMOTE-resampled training data (GBT-SMOTE).

\subsection{Imbalance Handling}

Three strategies were compared on LR to isolate resampling effects:
(1)~\texttt{class\_weight=`balanced'}; (2)~SMOTE~\cite{chawla2002smote}, which
interpolates between minority samples to achieve 1:1 training ratio; and
(3)~random under-sampling, which discards majority samples until 1:1.

\subsection{Threshold Tuning and Cost-Sensitive Evaluation}

For threshold $t$, predicted class is $\hat{y} = \mathbf{1}[\hat{p} \geq t]$.
We sweep $t \in [0.01, 0.99]$ (200 steps) on the validation set and report:

\begin{itemize}
    \item \textbf{F1-optimal threshold:} $t^*_{\text{F1}} = \arg\max_t \text{F1}(t)$
    \item \textbf{Cost-optimal threshold:}
          $t^*_{\text{cost}} = \arg\min_t \left[\text{FN}(t)\cdot C_{\text{FN}} + \text{FP}(t)\cdot C_{\text{FP}}\right]$
\end{itemize}

with $C_{\text{FN}} = 500$\,\euro{} (average fraud loss) and $C_{\text{FP}} = 10$\,\euro{}
(manual review cost per transaction). All threshold selection is done on the validation
set; the test set is evaluated exactly once.

% ─────────────────────────────────────────────────────────────────────────────
\section{Results}

\subsection{Resampling Strategy Comparison}

Table~\ref{tab:resampling} shows LR performance under the three imbalance strategies
at the default threshold of 0.50. Class weighting and SMOTE achieve similar PR-AUC;
random under-sampling underperforms due to the large information loss from discarding
199,012 majority-class training samples (99.8\% of legitimate transactions). SMOTE
is therefore applied to GBT as the secondary variant.

\begin{table}[ht]
    \centering
    \caption{Resampling strategy comparison on LR (validation set, threshold\,=\,0.50).}
    \label{tab:resampling}
    \begin{tabular}{lccc}
        \toprule
        \textbf{Strategy} & \textbf{PR-AUC} & \textbf{Recall} & \textbf{Fraud F1} \\
        \midrule
        Class weight   & 0.7520 & 0.9189 & 0.1265 \\
        SMOTE          & 0.7599 & 0.9324 & 0.1209 \\
        Under-sampling & 0.6986 & 0.9324 & 0.0865 \\
        \bottomrule
    \end{tabular}
\end{table}

\subsection{Model Comparison at F1-Optimal Threshold}

Table~\ref{tab:models} compares all model variants at their respective F1-optimal
thresholds on the validation set. RF achieves the highest PR-AUC (0.8637), confirming
that nonlinear feature interactions are important for fraud detection.
GBT-SMOTE achieves the highest F1 (0.8794) but at a lower PR-AUC (0.8533).
LR lags substantially on both PR-AUC and F1.

The F1-optimal thresholds vary widely: RF finds its optimum at 0.192 (well below
the default 0.5), while LR requires 0.990. This divergence reflects differences in
probability calibration across model families.

\begin{table}[ht]
    \centering
    \caption{All models at their F1-optimal thresholds (validation set). Best values in bold.}
    \label{tab:models}
    \setlength{\tabcolsep}{5pt}
    \begin{tabular}{lcccccc}
        \toprule
        \textbf{Model} & \textbf{PR-AUC} & \textbf{ROC-AUC}
            & \textbf{Thr.} & \textbf{F1} & \textbf{Recall} & \textbf{Prec.} \\
        \midrule
        LR (bal.)           & 0.7520 & 0.9860 & 0.990 & 0.7273 & 0.8649 & 0.6275 \\
        \textbf{RF (bal.)}  & \textbf{0.8637} & 0.9514 & \textbf{0.192}
                            & 0.8767 & \textbf{0.8649} & 0.8889 \\
        GBT (bal.)          & 0.7840 & \textbf{0.9822} & 0.951 & 0.8531 & 0.8243 & 0.8841 \\
        GBT (SMOTE)         & 0.8533 & 0.9729 & 0.960 & \textbf{0.8794} & 0.8378 & \textbf{0.9254} \\
        \bottomrule
    \end{tabular}
\end{table}

\subsection{Threshold Tuning and Cost Analysis}

Figure~\ref{fig:threshold} shows F1, Precision, and Recall as a function of threshold
for all models. Figure~\ref{fig:cost} shows total business cost vs.\ threshold.
Because $C_{\text{FN}} \gg C_{\text{FP}}$ (50:1), the cost-optimal thresholds are
substantially lower than the F1-optimal thresholds across all models---the system
prefers flagging more transactions (higher FP) to avoid missing fraud (lower FN).

For RF, the cost-optimal threshold drops to 0.044, catching one additional fraud
(FN: $13 \to 12$) while increasing FP from 11 to 48, for a net saving of \euro130
(\euro6{,}480 vs.\ \euro6{,}610). Table~\ref{tab:cost} summarises the comparison.

\begin{table}[ht]
    \centering
    \caption{F1-optimal vs.\ cost-optimal threshold analysis for all models on validation set.}
    \label{tab:cost}
    \setlength{\tabcolsep}{4pt}
    \begin{tabular}{lcc}
        \toprule
        \textbf{Model} & \textbf{F1-opt.\ Thr.} & \textbf{Cost-opt.\ Thr.} \\
        \midrule
        LR (balanced)  & 0.990 & 0.985 \\
        RF (balanced)  & 0.192 & 0.044 \\
        GBT (balanced) & 0.951 & 0.837 \\
        GBT (SMOTE)    & 0.960 & 0.251 \\
        \bottomrule
    \end{tabular}
\end{table}

\subsection{Final Test Set Evaluation}

RF (balanced) was selected as the best model by validation PR-AUC\,=\,0.8637.
Table~\ref{tab:test} reports test-set performance at both thresholds.

\begin{table}[ht]
    \centering
    \caption{RF (balanced) on the held-out test set (74 fraud; 42,648 legitimate).}
    \label{tab:test}
    \setlength{\tabcolsep}{5pt}
    \begin{tabular}{lcccccc}
        \toprule
        \textbf{Threshold} & \textbf{ROC-AUC} & \textbf{PR-AUC}
            & \textbf{F1} & \textbf{Recall} & \textbf{Prec.} & \textbf{Cost (\euro)} \\
        \midrule
        0.192 (F1-opt.)  & 0.9376 & 0.8116 & 0.8356 & 0.8243 & 0.8472 & 6{,}610 \\
        0.044 (cost-opt.)& 0.9376 & 0.8116 & 0.6739 & 0.8378 & 0.5636 & 6{,}480 \\
        \bottomrule
    \end{tabular}
\end{table}

At the F1-optimal threshold the model catches 61 of 74 frauds (82.4\% recall)
with only 11 false alarms. Figure~\ref{fig:test_eval} shows the confusion matrix,
ROC curve, and Precision-Recall curve on the test set.

% ─────────────────────────────────────────────────────────────────────────────
\section{Error Analysis}

The 13 false negatives (missed frauds at threshold 0.192) had a mean model score of
0.021, vs.\ 0.826 for detected frauds. Feature-distribution analysis of top
discriminative features (\texttt{V14}, \texttt{V17}, \texttt{V12}) shows that
missed frauds have values close to the legitimate distribution, indicating they are
genuinely ambiguous cases rather than model failures. Their transaction amounts were
slightly lower than detected frauds, consistent with small test charges being harder
to distinguish from low-value legitimate transactions. The 11 false positives had
a mean score of 0.44, just above the decision threshold, suggesting they are
legitimate transactions with unusual PCA projections resembling known fraud signatures.

% ─────────────────────────────────────────────────────────────────────────────
\section{Discussion and Conclusions}

\paragraph{Nonlinear models are essential.}
RF and GBT outperform LR by 0.11 PR-AUC, confirming that fraud signals arise from
nonlinear feature combinations that a linear boundary cannot capture.

\paragraph{Threshold tuning is as important as model choice.}
Dynamic tuning of the RF threshold from 0.50 to 0.192 improves F1 from 0.8613 to
0.8767 on the validation set with no model change.

\paragraph{Cost-sensitive analysis changes the operating point.}
The 50:1 FN/FP cost ratio shifts the RF optimal threshold from 0.192 to 0.044.
Whether this trade-off is acceptable depends on the institution's capacity for
manual fraud review.

\paragraph{SMOTE vs.\ class weighting.}
Both achieve similar PR-AUC on LR (+0.008 for SMOTE). On GBT, class weighting
yields higher PR-AUC while SMOTE yields higher F1; class weighting appears better
calibrated across all thresholds.

\paragraph{Limitations and future work.}
Features \texttt{V1}--\texttt{V28} are opaque PCA components; future work with raw
transaction features would enable actionable model explanations. Hyperparameter tuning
(grid search over tree depth, learning rate) and time-aware cross-validation (respecting
temporal ordering) remain as extensions. Online learning~\cite{pozzolo2014} could
address distribution shift as fraud strategies evolve.

% ── Figures ──────────────────────────────────────────────────────────────────

\begin{figure}[ht]
    \centering
    \includegraphics[width=0.95\textwidth]{figures/fig_11_threshold_tuning.png}
    \caption{F1-score, Precision, and Recall vs.\ decision threshold for all four models on
             the validation set. Dots mark the F1-maximising threshold per model.}
    \label{fig:threshold}
\end{figure}

\begin{figure}[ht]
    \centering
    \includegraphics[width=0.92\textwidth]{figures/fig_12_cost_analysis.png}
    \caption{Left: Total business cost vs.\ threshold for all models (validation set).
             Right: Comparison of F1-optimal vs.\ cost-optimal thresholds. Cost-optimal
             thresholds are consistently lower due to the 50:1 FN/FP cost ratio.}
    \label{fig:cost}
\end{figure}

\begin{figure}[ht]
    \centering
    \includegraphics[width=0.95\textwidth]{figures/fig_13_final_test_evaluation.png}
    \caption{Final test set evaluation for RF (balanced) at threshold 0.192.
             Confusion matrix (61 TP, 11 FP, 13 FN, 42,637 TN),
             ROC curve (AUC\,=\,0.9376), and PR curve (AUC\,=\,0.8116).}
    \label{fig:test_eval}
\end{figure}

% ─────────────────────────────────────────────────────────────────────────────
\begin{thebibliography}{9}

\bibitem{kaggle_ccfd}
    ULB Machine Learning Group,
    \textit{Credit Card Fraud Detection Dataset},
    Kaggle, 2018.
    \url{https://www.kaggle.com/datasets/mlg-ulb/creditcardfraud}

\bibitem{chawla2002smote}
    N.V. Chawla, K.W. Bowyer, L.O. Hall, W.P. Kegelmeyer,
    ``SMOTE: Synthetic Minority Over-sampling Technique,''
    \textit{J.\ Artificial Intelligence Research}, 16:321--357, 2002.

\bibitem{sklearn}
    F. Pedregosa et al.,
    ``Scikit-learn: Machine Learning in Python,''
    \textit{J.\ Machine Learning Research}, 12:2825--2830, 2011.

\bibitem{pozzolo2014}
    A. Dal Pozzolo, O. Caelen, R.A. Johnson, G. Bontempi,
    ``Calibrating Probability with Undersampling for Unbalanced Classification,''
    \textit{IEEE SSCI}, 2015.

\end{thebibliography}

\end{document}
